\begin{proposition}[Refuted: rational direct membership in P1]\label{B:P1-rational}
There is a small filtered category $A$ and a triple
$a_0,a_1,a_2:p\to q$ such that its rational column module
\[
 Q_a=\operatorname{coker}\bigl(
 \mathbb Q[A(q,-)]\longrightarrow\mathbb Q[A(p,-)]^2,\quad
 t\longmapsto(t(a_1-a_0),t(a_2-a_0))\bigr)
\]
does not belong to the smallest class of rational diagrams containing
all permutation modules and closed under extensions, kernels of
epimorphisms, arbitrary direct sums, direct summands, filtered colimits,
and tensoring with permutation modules. This strengthens the integral
torsion obstruction: all values of this example are rational vector
spaces. It does not refute the alternative in P1 of using another
module with the same $T_1$. Moreover, the same category and triple
give an integral column module with torsion-free values which lies
outside the original integral resolving class.
\end{proposition}
\begin{proof}
First fix an automorphism $\sigma:v\to v$ in a category. For a rational
diagram $M$, put $T_M=M(\sigma)$ and consider the condition
\begin{equation}\label{B:P1-rational-invariant}
 \ker(T_M^{2n}-2^{2n}\operatorname{id}_{M(v)})=0
 \qquad\text{for every integer }n\geq1.
\end{equation}
Every permutation module satisfies this condition. A finite-support
eigenvector of a permutation has support a finite union of finite
orbits. On an orbit of length $r$ its eigenvalue must have $r$th power
one; the eigenvalue $2^{2n}$ cannot occur.

Condition~\eqref{B:P1-rational-invariant} is preserved by subobjects,
extensions, arbitrary sums and summands. It is preserved by filtered
colimits because evaluation at $v$ and the displayed operator commute
with them, and filtered colimits of rational vector spaces are exact.
In particular it is preserved by kernels of epimorphisms.

We verify the additional tensor operation explicitly. Let $P$ be the
permutation of a set $S$ and give $M(v)\otimes\mathbb Q[S]$ the
automorphism $T_M\otimes P$. Fix $n\geq1$ and a finite-support vector
$z$ in the kernel of
$(T_M\otimes P)^{2n}-2^{2n}$. Decompose its support into $P^{2n}$-orbits.
On an infinite orbit, the eigenvector equation and invertibility of
$T_M$ propagate any nonzero coefficient indefinitely, contradicting
finite support. On an orbit of length $r$, iterating the coefficient
equation gives
\[
 (T_M^{2nr}-2^{2nr})m=0.
\]
The old condition forces $m=0$. Thus $z=0$, which proves tensor
closure. The use of all even powers is essential; injectivity of
$T_M+2$ alone would not suffice for this operation.

Now define $A$. Its objects are $p,q,v,t$, with $t$ terminal. The only
endomorphisms of $p,q,t$ are identities, and
$\operatorname{End}(v)=\{\sigma^n:n\in\mathbb Z\}$ is infinite cyclic.
The other nonterminal hom-sets are
\[
 A(p,q)=\{a_0,a_1,a_2\},\quad
 A(q,v)=\{h_{j,n}:j=1,2,\ n\in\mathbb Z\},\quad
 A(p,v)=\{b_n:n\in\mathbb Z\}.
\]
There are no further arrows except the unique arrows to $t$.
Composition with $\sigma^m$ adds $m$ to the subscript, and
\[
 h_{j,n}a_i=b_{n+e_{j,i}},\qquad
 (e_{1,0},e_{1,1},e_{1,2})=(0,1,2),\quad
 (e_{2,0},e_{2,1},e_{2,2})=(0,2,1).
\]
All composites with terminal arrows are the unique terminal arrow.
These formulas define an associative category: all nonidentity
nonterminal associativity checks are associativity of integer addition
or the equality $m+(n+e_{j,i})=(m+n)+e_{j,i}$; terminal composites are
unique. The object $t$ makes $A$ filtered.

Set $R=\mathbb Q[z,z^{-1}]$, using $z$ to represent the action of
$\sigma$. At $v$, the domain of the column presentation is $R^2$ and
each of its two target factors is $R$. The presentation matrix is
\[
 D=\begin{pmatrix}z-1&z^2-1\\z^2-1&z-1\end{pmatrix}.
\]
The sum and difference bases diagonalize it over $R$, giving entries
$(z-1)(z+2)$ and $-z(z-1)$. Since $-z$ is a unit,
\[
 Q_a(v)\cong R/((z-1)(z+2))\oplus R/(z-1).
\]
For an explicit nonzero $-2$ eigenvector, let
$\xi$ be the class of $(z-1,z-1)$. Then
\[
 (z+2)(z-1,z-1)=D(1,1),
\]
so $(T_{Q_a}+2)\xi=0$ and consequently
$(T_{Q_a}^2-4)\xi=0$. The functional obtained by evaluating at $z=-2$
and summing the two coordinates annihilates both columns of $D$ but
takes value $-6$ on $(z-1,z-1)$. Thus $\xi\ne0$.

The column module fails~\eqref{B:P1-rational-invariant} for $n=1$,
whereas the entire specified resolving class satisfies it. Hence it
does not belong to that class. Since $t$ is terminal,
$\colim_A Q_a=Q_a(t)=\mathbb Q^2$ and $T_i(Q_a)=0$ for $i>0$.
This example concerns direct membership only, not the main implication.

\paragraph{Proved integral strengthening.}
The same invariant and closure argument work over $\mathbb Z$.
For the tensor operation on an infinite permutation orbit, choose an
extreme nonzero coefficient in its finite support. The coefficient
equation one step beyond that support says that $T_M^{2n}$ kills
that coefficient, contrary to invertibility of $T_M$. This argument
does not divide by $2$; the finite-orbit argument is unchanged.

Let $R_{\mathbb Z}=\mathbb Z[z,z^{-1}]$. The integral column matrix at
$v$ is the same $D$, and
\[
 \det D=-z(z-1)^2(z+2)=-z^4+3z^2-2z.
\]
It is nonzero over the domain $R_{\mathbb Z}$, so
$0\to R_{\mathbb Z}^2\xrightarrow{D}R_{\mathbb Z}^2
\to Q_a(v)\to0$ is exact. Its reduction modulo any prime $p$
is also injective: the determinant has leading coefficient $-1$,
so it remains nonzero in the domain $\mathbb F_p[z,z^{-1}]$.
The two free abelian terms in this presentation therefore give
\[
 \operatorname{Tor}^{\mathbb Z}_1(\mathbb Z/p,Q_a(v))
 =\ker(D\bmod p)=0
 \qquad\text{for every prime }p.
\]
Equivalently, $Q_a(v)$ has no element killed by a prime and is
torsion free. At $p$ and $t$ the column module is $\mathbb Z^2$;
at $q$ it is $\mathbb Z^5$, since the single presentation column
has a coordinate equal to $1$. Thus every value is torsion free.
The same vector $\xi=[(z-1,z-1)]$ is killed by $z+2$ and is
nonzero: evaluate at $z=-2$ in $\mathbb Z[1/2]$, then sum the
coordinates, obtaining $-6$. It violates the integral invariant
for $n=1$. Hence the integral column module also fails membership
in P1's full resolving class, although all its higher colimits
vanish because $t$ is terminal.
\end{proof}

\paragraph{Exact symbolic certificate.}
The program \texttt{finite\_rational\_p1.py} checks every composable
triple of the finitely many arrow families using three independent
integer exponent variables. Equality of the resulting affine
expressions proves the displayed associativity identities for all
integer exponents; no finite truncation of the category is substituted.
It also verifies the rational change of basis, the two polynomial
relations for $\xi$, and the nonzero detecting functional. The full
tensor-closure argument is the orbit proof above, not a numerical test.
